\chapter{Banco de Dados}
\label{cap:banco-dados}

\begin{edital}
\textbf{Atenção ao recorte:} o Perfil 3 \emph{não tem} ``Banco de Dados''
como disciplina. O assunto entra por duas portas: \textbf{Inteligência de
Negócios} (``arquitetura e aplicações de \emph{data warehouse} com ETL e
OLAP'', ``\emph{data warehouse} e \emph{data mining}'', ``visualização de
dados: BD individuais e cubos'') e o item 20 de Desenvolvimento de Sistemas
(``Conceitos de Inteligência Artificial, \textbf{Análise de Dados} e Big
Data'').

A lista fechada --- modelagem e normalização, arquitetura, estrutura de
dados, SQL (ANSI), administração de dados, backup/restauração, engenharia de
dados/Big Data --- e os \textbf{SGBD nomeados} (Oracle 19C, MySQL,
PostgreSQL, MongoDB, MS-SQL Server 2019) estão no \textbf{Perfil 2}. Estudamos
assim mesmo porque a \textbf{FGV cobrou banco de dados no Perfil 3 da Dataprev
2024} --- do mesmo jeito que cobrou redes, que também não está no perfil.
\end{edital}
\peso{ALTO \textbf{por evidência de prova}, não por peso de edital --- metade
do ``eixo duplo'' com Engenharia de Software; no MPU 2025 foi o bloco que mais
caiu.}

\section{Modelagem em três níveis}

\begin{center}
\begin{tabular}{@{}p{2.4cm}p{4.6cm}p{4.6cm}@{}}
\toprule
\textbf{Nível} & \textbf{O que é} & \textbf{Artefato} \\
\midrule
Conceitual & visão de negócio, independe de SGBD & Modelo Entidade-Relacionamento (MER/DER) \\
Lógico & estrutura no modelo escolhido (relacional) & tabelas, chaves, tipos genéricos \\
Físico & implementação no SGBD específico & DDL, índices, particionamento \\
\bottomrule
\end{tabular}
\end{center}

\begin{pegadinha}
Ordem correta: conceitual $\to$ lógico $\to$ físico. A FGV inverte a ordem ou
troca o artefato de cada nível.
\end{pegadinha}

\section{Modelo relacional e integridade}

\begin{itemize}
\item \textbf{Chave primária (PK):} identifica unicamente a tupla; não nula,
única.
\item \textbf{Chave estrangeira (FK):} referencia a PK de outra (ou da
mesma) tabela.
\item \textbf{Integridade referencial:} toda FK aponta para uma PK
existente (ou é nula). Ações: CASCADE, SET NULL, RESTRICT/NO ACTION.
\item \textbf{Integridade de entidade:} PK não pode ser nula.
\item \textbf{Álgebra relacional:} seleção ($\sigma$) = \texttt{WHERE} =
filtra \textbf{linhas}; projeção ($\pi$) = filtra \textbf{colunas}.
\end{itemize}

\subsection{Do MER para o relacional (como cada cardinalidade vira tabela)}

\begin{center}
\begin{tabular}{@{}p{2.2cm}p{9.8cm}@{}}
\toprule
\textbf{Cardinalidade} & \textbf{Implementação no modelo relacional} \\
\midrule
1:1 & FK em uma das duas tabelas (de preferência no lado de participação obrigatória), com restrição de unicidade \\
1:N & \textbf{FK no lado N} --- a tabela ``muitos'' guarda a chave da tabela ``um'' \\
\textbf{N:M} & \textbf{tabela associativa} (intermediária), com as FKs das duas pontas formando a chave primária composta \\
\bottomrule
\end{tabular}
\end{center}

\begin{pegadinha}
N:M \textbf{nunca} se resolve com FK direta em uma das tabelas (isso só
atende 1:N) nem com coluna multivalorada dos dois lados --- coluna
multivalorada viola a \textbf{1FN}. Trigger e índice composto não implementam
cardinalidade nenhuma: são distratores de outra prateleira.
\end{pegadinha}

\begin{conceito}[Entidade fraca $\times$ entidade associativa]
\textbf{Entidade fraca:} não se identifica sozinha. Sua chave é a chave da
\textbf{entidade proprietária} (forte) somada a uma \textbf{chave parcial}
(discriminador) --- ex.: DEPENDENTE só é identificado dentro do cadastro de um
SERVIDOR, pelo nome. O relacionamento com o proprietário é
\textbf{identificador}, e a existência da fraca depende da forte.

\textbf{Entidade associativa:} é outra coisa --- é o \textbf{relacionamento
N:M promovido a entidade}, quando ele próprio precisa se relacionar com uma
terceira entidade ou ganhar atributos próprios.
\end{conceito}

\begin{pegadinha}
A quase-certa diz ``entidade fraca, que \textbf{dispensa} a chave do
proprietário'' --- é justamente o oposto: sem a chave do proprietário ela não
se identifica. E ter atributos próprios \textbf{não} torna a entidade forte; o
que define a força é conseguir se identificar sozinha.
\end{pegadinha}

\section{Normalização (formas normais)}

\begin{center}
\begin{tabular}{@{}p{2cm}p{9.4cm}@{}}
\toprule
\textbf{FN} & \textbf{Elimina} \\
\midrule
1FN & atributos multivalorados / não atômicos (cada célula um valor) \\
2FN & dependência parcial da chave (só relevante com PK composta) \\
3FN & dependência transitiva (atributo que depende de outro não-chave) \\
FNBC (BCNF) & anomalia quando há mais de uma chave candidata sobreposta \\
\bottomrule
\end{tabular}
\end{center}

Ordem: 1FN $\to$ 2FN $\to$ 3FN $\to$ BCNF. Normalizar reduz redundância;
\textbf{desnormalizar de propósito} (ex.: Data Warehouse) é aceitável por
desempenho de leitura --- objetivo é leitura, nunca escrita.

\section{SQL: DDL $\times$ DML $\times$ DCL $\times$ TCL}

\begin{center}
\begin{tabular}{@{}p{2.6cm}p{4cm}p{4.5cm}@{}}
\toprule
\textbf{Sublinguagem} & \textbf{Comandos} & \textbf{Observação} \\
\midrule
DDL (definição) & CREATE, ALTER, \textbf{DROP, TRUNCATE} & TRUNCATE é DDL (não DML), auto-commit \\
DML (manipulação) & SELECT, INSERT, UPDATE, DELETE & DELETE é DML e pode ter WHERE + rollback \\
DCL (controle) & GRANT, REVOKE & permissões \\
TCL (transação) & COMMIT, ROLLBACK, SAVEPOINT & \\
\bottomrule
\end{tabular}
\end{center}

\begin{pegadinha}
\textbf{TRUNCATE (DDL) $\times$ DELETE (DML)}; \texttt{UNION} remove
duplicatas / \texttt{UNION ALL} mantém; \texttt{WHERE} filtra linhas /
\texttt{HAVING} filtra grupos após \texttt{GROUP BY}; \texttt{JOIN} sem
\texttt{ON} vira produto cartesiano. Ordem de execução lógica: FROM $\to$
WHERE $\to$ GROUP BY $\to$ HAVING $\to$ SELECT $\to$ ORDER BY. Leia palavra
por palavra os comandos SQL da questão --- a FGV troca uma cláusula só
(LIMIT/OFFSET, DISTINCT, ILIKE, NOT EXISTS por JOIN comum) e testa leitura de
código.
\end{pegadinha}

\subsection{VIEW, GRANT e o controle da transação}

\textbf{VIEW (visão):} uma \textbf{tabela virtual} --- é a consulta guardada
com nome, não os dados copiados. Consultar a visão executa o \texttt{SELECT}
que está por trás dela, então o resultado é \textbf{sempre atual}. Serve para
simplificar consulta complexa e para \textbf{restringir acesso} (a visão expõe
só as colunas/linhas permitidas, e o \texttt{GRANT} é dado sobre ela em vez da
tabela). Visão \textbf{simples} costuma aceitar atualização; visão com
\texttt{JOIN}, agregação ou \texttt{DISTINCT}, não. A
\textbf{\emph{materialized view}} é a exceção que confirma a regra: essa
\emph{sim} armazena o resultado em disco e precisa ser atualizada.

\textbf{GRANT/REVOKE (DCL):} \texttt{GRANT SELECT, INSERT ON tabela TO
usuario} concede; \texttt{REVOKE} retira. Com \texttt{WITH GRANT OPTION}, quem
recebeu pode repassar o privilégio a terceiros --- detalhe que a banca gosta
de cobrar. \textbf{GRANT é DCL}, não DDL nem DML.

\textbf{SAVEPOINT e ROLLBACK TO (TCL):} o \texttt{SAVEPOINT nome} marca um
ponto \textbf{dentro} da transação; \texttt{ROLLBACK TO nome} desfaz só o que
veio \textbf{depois} daquela marca --- e a transação \textbf{continua aberta},
sem \texttt{COMMIT} nem \texttt{ROLLBACK} total. É desfazer parcial.

\begin{pegadinha}
Três trocas prontas: dizer que a \textbf{visão armazena os dados}
fisicamente (isso é a \emph{materialized view}; a visão comum é virtual);
classificar \textbf{GRANT como DDL} (é \textbf{DCL}); e afirmar que o
\texttt{ROLLBACK TO} \textbf{encerra} a transação ou desfaz tudo --- ele
volta ao \emph{savepoint} e a transação segue viva.
\end{pegadinha}

\subsection{Notações do MER e chave surrogada}

\textbf{Crow's Foot (pé de galinha):} a notação de cardinalidade mais usada em
ferramenta. Lê-se \textbf{no ponto onde a linha encosta na entidade}, e cada
lado traz \emph{dois} símbolos --- o de fora é o \textbf{máximo}, o de dentro
é o \textbf{mínimo}:

\begin{center}
\begin{tabular}{@{}p{4.6cm}p{7.4cm}@{}}
\toprule
\textbf{Símbolo} & \textbf{Leitura} \\
\midrule
Traço (\emph{barra}) & exatamente \textbf{um} (mínimo 1 / máximo 1) \\
Círculo (\emph{o}) & \textbf{zero} --- opcional \\
Pé de galinha (três riscos) & \textbf{muitos} \\
Círculo + pé de galinha & \textbf{zero ou muitos} \\
Traço + pé de galinha & \textbf{um ou muitos} (obrigatório) \\
\bottomrule
\end{tabular}
\end{center}

Daí a tradução para DDL: o lado com \textbf{pé de galinha} é o lado
\textbf{N}, e é nele que entra a \textbf{FK}; o círculo indica participação
\textbf{opcional} $\to$ a FK aceita \texttt{NULL}; o traço indica participação
\textbf{obrigatória} $\to$ \texttt{NOT NULL}.

\textbf{Chave surrogada (substituta/artificial):} chave primária
\textbf{sem significado de negócio}, gerada pelo sistema (sequência,
\emph{identity}, UUID), em oposição à \textbf{chave natural}, que vem do
domínio (CPF, matrícula). Vantagens: é estável (o dado de negócio pode mudar
--- CPF corrigido, matrícula reemitida), curta e uniforme para \emph{join}. É
padrão em \textbf{Data Warehouse}, onde ela também permite guardar
\textbf{versões históricas} da mesma entidade de negócio (dimensão de
mudança lenta) --- detalhe no Capítulo \ref{cap:bi}.

\begin{pegadinha}
Chave surrogada \textbf{não é} chave estrangeira nem chave composta, e não
dispensa a chave natural: o campo de negócio continua na tabela, em geral com
restrição \texttt{UNIQUE}. Em Crow's Foot, a inversão clássica é ler a
cardinalidade \textbf{do lado errado} da linha --- o símbolo vale para a
entidade que ele \emph{toca}.
\end{pegadinha}

\subsection{Código no servidor: gatilho $\times$ procedimento armazenado}

\begin{center}
\begin{tabular}{@{}p{2.4cm}p{4.8cm}p{4.8cm}@{}}
\toprule
& \textbf{Gatilho (trigger)} & \textbf{Procedimento armazenado} \\
\midrule
Quem dispara & o \textbf{próprio SGBD}, em resposta a um \textbf{evento} de dados & a \textbf{aplicação} (ou outro código), por chamada explícita \\
Como se invoca & não se invoca --- é automático & \texttt{CALL}/\texttt{EXECUTE} \\
Amarrado a & uma tabela + evento (INSERT, UPDATE, DELETE) & nada; é código nomeado e reutilizável \\
Parâmetros & não recebe & recebe (e pode retornar) \\
Bom para & auditoria, log, regra que \textbf{não pode depender} da aplicação lembrar & rotina de negócio chamada sob demanda \\
\bottomrule
\end{tabular}
\end{center}

\begin{pegadinha}
A troca é sempre a mesma: dizer que o \textbf{procedimento} é ``acionado
automaticamente pelo SGBD a cada UPDATE'' (isso é o gatilho) ou que o
\textbf{gatilho} é ``invocado com um \texttt{CALL} antes do UPDATE'' (isso é o
procedimento). O gatilho \textbf{também não substitui} restrição de
integridade referencial declarada --- são mecanismos diferentes. Palavra-chave
do enunciado: ``\textbf{sem depender de a aplicação lembrar}'' $\to$ gatilho.
\end{pegadinha}

\section{Propriedades ACID (transações)}

\begin{itemize}
\item \textbf{A}tomicidade: tudo ou nada.
\item \textbf{C}onsistência: leva o banco de um estado válido a outro.
\item \textbf{I}solamento: transações concorrentes não interferem.
\item \textbf{D}urabilidade: efeito confirmado persiste (mesmo após falha).
\end{itemize}

\section{Concorrência: níveis de isolamento}

O ``I'' do ACID não é tudo-ou-nada: o padrão SQL define \textbf{quatro
níveis}, e cada um \emph{admite} certos fenômenos de concorrência. Quanto
maior o isolamento, \textbf{menor} a concorrência --- é um trade-off, não um
ganho de graça.

\textbf{Os três fenômenos:}
\begin{itemize}
\item \textbf{Leitura suja} (\emph{dirty read}): lê dado alterado por outra
transação que \textbf{ainda não confirmou} (e pode desfazer).
\item \textbf{Leitura não repetível} (\emph{non-repeatable read}): relê a
\textbf{mesma linha} e o valor mudou, porque outra transação a atualizou e
confirmou no meio.
\item \textbf{Leitura fantasma} (\emph{phantom read}): relê a \textbf{mesma
faixa} e aparecem \textbf{linhas novas}, inseridas e confirmadas por outra
transação.
\end{itemize}

\begin{center}
\begin{tabular}{@{}p{3.6cm}p{2.6cm}p{2.6cm}p{2.6cm}@{}}
\toprule
\textbf{Nível} & \textbf{Leitura suja} & \textbf{Não repetível} & \textbf{Fantasma} \\
\midrule
READ UNCOMMITTED & admite & admite & admite \\
READ COMMITTED & \textbf{impede} & admite & admite \\
REPEATABLE READ & impede & \textbf{impede} & admite \\
SERIALIZABLE & impede & impede & \textbf{impede} \\
\bottomrule
\end{tabular}
\end{center}

\begin{pegadinha}
Duas inversões prontas: dizer que o \textbf{READ UNCOMMITTED é o mais
restritivo} (é o \textbf{mais permissivo} --- é justamente ele que deixa ler
o não confirmado) e dizer que o \textbf{REPEATABLE READ elimina o fantasma}
(elimina a \textbf{não repetível}; o fantasma só cai no SERIALIZABLE). E o
absoluto invertido: ``quanto mais alto o isolamento, maior a concorrência''
--- é o contrário. Marque o par pelo objeto: linha que \textbf{muda de valor}
= não repetível; linha \textbf{nova que aparece} = fantasma.
\end{pegadinha}

\section{Relacional $\times$ Multidimensional (OLTP $\times$ OLAP)}

\begin{center}
\begin{tabular}{@{}p{2.6cm}p{4.4cm}p{4.4cm}@{}}
\toprule
& \textbf{OLTP (transacional)} & \textbf{OLAP (analítico)} \\
\midrule
Uso & operações do dia a dia & apoio à decisão / análise \\
Escrita/leitura & muitas escritas curtas & leitura pesada, agregação \\
Modelo & relacional normalizado & multidimensional (cubos) \\
Estrutura & tabelas & dimensões e métricas/fatos \\
\bottomrule
\end{tabular}
\end{center}

\begin{pegadinha}
Pegadinha central do bloco: a FGV troca OLTP$\leftrightarrow$OLAP. ``Cubos,
dimensões, agregação, análise'' = OLAP; ``transação, muitas gravações'' =
OLTP. Nunca ``relacional é usado em OLAP e multidimensional em OLTP''.
\end{pegadinha}

\section{NoSQL e novos armazenamentos}

\begin{itemize}
\item \textbf{NoSQL:} alta disponibilidade e \textbf{escalabilidade
horizontal}; costuma relaxar ACID (modelo \textbf{BASE}: Basically
Available, Soft state, Eventual consistency). Tipos: \textbf{chave-valor}
(Redis), \textbf{documento} (MongoDB), \textbf{colunar} (Cassandra),
\textbf{grafo} (Neo4j, nós e arestas). Não é ``sempre grafo'' nem ``segue
ACID estritamente''; não substitui ERP/CRM transacional por padrão.
\item \textbf{Teorema CAP:} em sistema distribuído, escolha 2 de 3 ---
\textbf{C}onsistency, \textbf{A}vailability, \textbf{P}artition tolerance.
Com partição (P inevitável em rede), decide-se entre C e A. Cenário de
registro oficial/consistência crítica (sentença, transação) $\to$ CP.
\item \textbf{Banco em memória} (in-memory, ex.: Redis): baixa latência.
\item \textbf{Data Lake $\times$ Data Warehouse $\times$ Lakehouse:} Lake
guarda dado bruto (schema-on-read, estruturado e não estruturado); DW guarda
dado tratado e modelado (schema-on-write); Lakehouse combina os dois.
\end{itemize}

\begin{pegadinha}
Absolutos denunciam o distrator: ``NoSQL segue ESTRITAMENTE ACID'', ``NoSQL
usa SEMPRE modelo relacional'', ``dados SEMPRE em grafos'', ``desnormalização
é requisito OBRIGATÓRIO''. A FGV adora ``sempre/apenas/nunca/estritamente'' ---
quase sempre é a alternativa errada. No teorema CAP, desconfie de cenários que
priorizam disponibilidade/velocidade (AP) quando o contexto pede consistência
crítica (deveria ser CP).
\end{pegadinha}

\section{ETL $\times$ ELT}

\begin{center}
\begin{tabular}{@{}p{2.6cm}p{4.3cm}p{4.3cm}@{}}
\toprule
& \textbf{ETL} & \textbf{ELT} \\
\midrule
Ordem & Extrai $\to$ \textbf{Transforma} $\to$ Carrega & Extrai $\to$ Carrega $\to$ \textbf{Transforma} \\
Transformação & antes de carregar (fora do destino) & depois, no próprio destino \\
Bom para & DW tradicional & destino com muito poder de processamento (nuvem, big data) \\
\bottomrule
\end{tabular}
\end{center}

\begin{pegadinha}
A pegadinha é a \textbf{posição do T}. ELT aproveita o processamento do
destino (bom para grande volume); ETL é melhor quando a transformação ocorre
antes/fora. Cuidado com a alternativa que inverte quem transforma antes/depois
do carregamento, ou que diz que ELT é melhor para volume pequeno (é o
contrário).
\end{pegadinha}

\section{Índices e notações}

\begin{itemize}
\item \textbf{Índices:} \textbf{B+ Tree} (padrão, bom para faixas e alta
cardinalidade), \textbf{bitmap} (baixa cardinalidade, ex.: sexo/status),
\textbf{hash} (igualdade exata). \textbf{Hashing extensível:} hash dinâmico
com um \textbf{diretório} de ponteiros e \textit{global depth}; quando um
bucket enche, ele faz \textbf{split} e dobra o diretório se preciso --- cresce
sem reorganizar tudo.
\item \textbf{IDEF1X:} notação de modelagem de dados (comum em ferramentas
como o erwin). Relacionamento \textbf{identificador} = linha \textbf{sólida};
entidade dependente = retângulo de \textbf{cantos arredondados}.
\item Administração física (tablespaces, ARIES/recuperação, backup,
otimizador) está no Capítulo \ref{cap:orfaos} (Órfãos).
\end{itemize}

\begin{jacaiu}
\textbf{Em prova real da FGV:} quase toda a lista deste bloco tem lastro.
Relacional $\times$ multidimensional (OLTP/OLAP); NoSQL; ETL $\times$ ELT
(``assinale a \textbf{incorreta}''); ETL como ponte para o DW ---
\textbf{Dataprev 2024}. Álgebra relacional (a seleção $\sigma$ e sua tradução
para o \texttt{WHERE}); \textbf{VIEW} como tabela virtual; \textbf{GRANT} de
privilégios; \textbf{SAVEPOINT} com \texttt{ROLLBACK TO} e
\texttt{RELEASE}; \textbf{Crow's Foot} $\to$ DDL, que no mesmo item cobra
\textbf{N:M por tabela associativa} com chave composta e \textbf{integridade
referencial}; \texttt{DROP} $\times$ \texttt{TRUNCATE} $\times$
\texttt{DELETE}; \texttt{LIKE}; MongoDB (relação $\leftrightarrow$ coleção,
\texttt{\$size}, ETL de NoSQL para relacional); NoSQL de \textbf{grafos} para
rede de relacionamentos; DAMA-DMBOK e governança de dados --- \textbf{MPU}.
Teorema \textbf{CAP} (o perfil CP); desnormalização intencional no DW; Star
$\times$ Snowflake; \textbf{Data Warehouse $\times$ Data Lake $\times$
Lakehouse}; \textbf{chave surrogada}; \texttt{LIMIT}/\texttt{OFFSET};
\texttt{DISTINCT}; anti-join (Partes sem Audiência); \emph{constraints}
(\texttt{UNIQUE}, \texttt{NOT NULL}, \texttt{CHECK}, PK autogerada) ---
\textbf{TJ-RJ}. \textbf{Entidade fraca} na notação IDEF1X; integridade
referencial; normalização (chave candidata e a forma normal mais alta, a
partir das dependências funcionais); \textbf{gatilho $\times$ procedimento
$\times$ função}, em três questões distintas; VIEW para esconder a coluna de
salário; NoSQL chave-valor para cache de sessão --- \textbf{ALERO 2026}.

\textbf{E tem mais 34, que só apareceram no fechamento do ITEM 7:} o bloco de
\textbf{administração física} da \textbf{ALERO 2026} vinha classificado como
``órfão'' por um defeito do importador. São \textbf{B+ Tree} (eficiência e
fator de ramificação) e \textbf{hashing extensível}; \textbf{índice
\emph{clustered}}; \textbf{ARIES/WAL} e a ordem das três fases;
\textbf{otimizador} de consultas e estatísticas desatualizadas;
\textbf{tablespaces}; \textbf{\emph{connection pooling}}; \textbf{XML Schema}
e XQuery no SGBD; \textbf{backup} com RPO/RTO; \textbf{Administração de Dados
$\times$ DBA}; \textbf{JDBC/ODBC}; \textbf{dicionário de dados}; e
\textbf{sharding}. Com elas, banco de dados salta de 25 para \textbf{59}
questões reais. O conteúdo de administração física continua sendo ensinado no
Capítulo \ref{cap:orfaos}.

\textbf{Não leia esse 59 como peso de prova.} \textbf{28 das 59} vêm de uma
prova só --- a da ALERO para o perfil \emph{Banco de Dados}, que não é o nosso.
Nas provas mais próximas de desenvolvimento, banco de dados fica em
\textbf{30}, atrás de engenharia de software (39). O número grande mede
\emph{material de treino disponível}, não a chance de cair: na Dataprev 2024,
banco de dados e BI somaram \textbf{4} questões. Estude como pilar --- é o
eixo duplo com engenharia de software ---, mas sem inflar.

\textbf{No nosso banco} (previsto pelo edital, ainda não visto na amostra de
provas): \textbf{níveis de isolamento} e os fenômenos que cada um admite
(leitura suja, não repetível, fantasma) --- nenhuma das 432 questões reais
menciona \texttt{READ COMMITTED} ou \texttt{SERIALIZABLE}.

Rode \texttt{./quiz.py banco-dados} e \texttt{./quiz.py bi}.
\end{jacaiu}

\section{Alta probabilidade / pesquisa extra}

\begin{itemize}
\item \textbf{SGBD nomeados no edital:} Oracle 19C, MySQL, PostgreSQL,
MongoDB, MS-SQL Server 2019 --- item \textbf{7.1 do Perfil 2}, não do 3.
Nenhuma questão deve depender de sintaxe proprietária de um deles; o que pode
vir para nós é o \textbf{conceito} (SQL ANSI, o par relacional $\times$
documento).
\item \textbf{Modelagem dimensional (Kimball):} \textbf{Star Schema} (uma
tabela fato + dimensões desnormalizadas, mais rápido) $\times$
\textbf{Snowflake} (dimensões normalizadas, menos redundância, mais joins)
--- detalhe completo no Capítulo \ref{cap:bi} (BI).
\item \textbf{Big Data --- os Vs.} O trio \textbf{clássico} (Gartner, 2001) é
\textbf{Volume, Velocidade e Variedade}. \textbf{Veracidade} e \textbf{Valor}
são \textbf{extensões} posteriores (daí os modelos ``4V'' e ``5V''), assim
como \textbf{Variabilidade}. Saiba qual é o trio original: a pegadinha é
trocar um membro do trio por um V de extensão --- ``volume, velocidade e
valor'' ou ``volume, variedade e veracidade'' são as ``quase certas''.
\item \textbf{MongoDB:} query, \texttt{\$size}, coleção/documento (vs.\
tabela/linha do MySQL); \textbf{banco orientado a grafos} tem nós e arestas.
\end{itemize}

\begin{comosair}
Memorize lado a lado: OLTP = transacional, normalizado, escrita, linha; OLAP
= analítico, dimensional, leitura, cubo/agregação. ETL = transforma ANTES (DW
tradicional); ELT = carrega e transforma no destino (nuvem, grande volume).
Desnormalização = menos joins, mais espaço, risco de anomalia de atualização
--- serve à leitura, não à escrita. NoSQL: BASE, alta disponibilidade e
escalabilidade horizontal; ruim para ERP/CRM fortemente relacional.
\end{comosair}
